\QHeading{第1题　在$\bigcirc$和$\square$里分别填写相应的运算符号和数。}
\par\noindent


\Equation{49\times 72+72\times 51=\textcolor{answerred}{72}\textcolor{answerred}{\times}(\textcolor{answerred}{49}\textcolor{answerred}{+}\textcolor{answerred}{51})}
\par

\Equation{(18+25)\times 4=\textcolor{answerred}{18}\times\textcolor{answerred}{4}\textcolor{answerred}{+}\textcolor{answerred}{25}\times\textcolor{answerred}{4}}
\Equation{39\times 99+39=\textcolor{answerred}{39}\times (\textcolor{answerred}{99}\textcolor{answerred}{+}\textcolor{answerred}{1})}
\Equation{65\times 15+65\times 85=(\textcolor{answerred}{15}\textcolor{answerred}{+}\textcolor{answerred}{85})\times\textcolor{answerred}{65}}
\par\noindent

\QHeading{第2题　用简便方法计算下面各题。}

\CalcBlock[方法一：简便计算]{1}{\begin{aligned}[t]&\phantom{={}}\color{black}(80-8)\times 125\\&\color{answerred}=80\times125-8\times125\\&\color{answerred}=10000-1000\\&\color{answerred}=9000\end{aligned}}
\par
\MethodBlock{方法二：先处理括号或乘除，再拆因数，凑整计算}{把一个因数拆成两个因数，先凑成整百或整千。}{\begin{aligned}[t]&\phantom{={}}(80-8)\times 125\\
&=72\times 125\\
&=125\times (8\times 9)\\
&=(125\times 8)\times 9\\
&=1000\times 9\\
&=9000\end{aligned}}

\MethodBlock{方法三：先处理括号或乘除，再把特殊因数改成整百或整千}{先把25、125等改写成整百或整千除以一个数。}{\begin{aligned}[t]&\phantom{={}}(80-8)\times 125\\
&=72\times 125\\
&=(1000\div 8)\times 72\\
&=1000\times (72\div 8)\\
&=1000\times 9\\
&=9000\end{aligned}}

\MethodBlock{方法四：先处理括号或乘除，再凑整后补偿}{把接近整十、整百的因数凑整，再补回或减去多算的部分。}{\begin{aligned}[t]&\phantom{={}}(80-8)\times 125\\
&=72\times 125\\
&=125\times (70+2)\\
&=125\times 70+125\times 2\\
&=8750+250\\
&=9000\end{aligned}}


\CalcBlock[方法一：简便计算]{2}{\begin{aligned}[t]&\phantom{={}}\color{black}175\times 18-75\times 18\\&\color{answerred}=(175-75)\times18\\&\color{answerred}=100\times18\\&\color{answerred}=1800\end{aligned}}
\MethodBlock{方法二：按运算顺序直接计算（对照）}{本法只作结果核对；本题要求简便计算，前面的简便解法仍保留。}{\begin{aligned}[t]&\phantom{={}}175\times 18-75\times 18\\
&=3150-1350\\
&=1800\end{aligned}}

\CalcBlock[方法一：简便计算]{3}{\begin{aligned}[t]&\phantom{={}}\color{black}32\times 99\\&\color{answerred}=32\times(100-1)\\&\color{answerred}=32\times100-32\times1\\&\color{answerred}=3200-32\\&\color{answerred}=3168\end{aligned}}
\MethodBlock{方法二：凑整后补偿}{把接近整十、整百的因数凑整，再补回或减去多算的部分。}{\begin{aligned}[t]&\phantom{={}}32\times 99\\
&=99\times (30+2)\\
&=99\times 30+99\times 2\\
&=2970+198\\
&=3168\end{aligned}}

\MethodBlock{方法三：按数位拆分计算}{把一个因数拆成整十、整百数与余下部分之和。}{\begin{aligned}[t]&\phantom{={}}32\times 99\\
&=32\times (90+9)\\
&=32\times 90+32\times 9\\
&=2880+288\\
&=3168\end{aligned}}

\CalcBlock[方法一：简便计算]{4}{\begin{aligned}[t]&\phantom{={}}\color{black}56\times 99+56\\&\color{answerred}=56\times(99+1)\\&\color{answerred}=56\times100\\&\color{answerred}=5600\end{aligned}}
\MethodBlock{方法二：按运算顺序直接计算（对照）}{本法只作结果核对；本题要求简便计算，前面的简便解法仍保留。}{\begin{aligned}[t]&\phantom{={}}56\times 99+56\\
&=5544+56\\
&=5600\end{aligned}}

\CalcBlock[方法一：简便计算]{5}{\begin{aligned}[t]&\phantom{={}}\color{black}125\times 81-125\\&\color{answerred}=125\times(81-1)\\&\color{answerred}=125\times80\\&\color{answerred}=10000\end{aligned}}
\MethodBlock{方法二：按运算顺序直接计算（对照）}{本法只作结果核对；本题要求简便计算，前面的简便解法仍保留。}{\begin{aligned}[t]&\phantom{={}}125\times 81-125\\
&=10125-125\\
&=10000\end{aligned}}

\CalcBlock[方法一：简便计算]{6}{\begin{aligned}[t]&\phantom{={}}\color{black}45\times 102\\&\color{answerred}=45\times(100+2)\\&\color{answerred}=45\times100+45\times2\\&\color{answerred}=4500+90\\&\color{answerred}=4590\end{aligned}}
\MethodBlock{方法二：按数位拆分计算}{把一个因数拆成整十、整百数与余下部分之和。}{\begin{aligned}[t]&\phantom{={}}45\times 102\\
&=102\times (40+5)\\
&=102\times 40+102\times 5\\
&=4080+510\\
&=4590\end{aligned}}

\CalcBlock[方法一：简便计算]{7}{\begin{aligned}[t]&\phantom{={}}\color{black}83\times 101-83\\&\color{answerred}=83\times(101-1)\\&\color{answerred}=83\times100\\&\color{answerred}=8300\end{aligned}}
\MethodBlock{方法二：按运算顺序直接计算（对照）}{本法只作结果核对；本题要求简便计算，前面的简便解法仍保留。}{\begin{aligned}[t]&\phantom{={}}83\times 101-83\\
&=8383-83\\
&=8300\end{aligned}}

\CalcBlock[方法一：简便计算]{8}{\begin{aligned}[t]&\phantom{={}}\color{black}25\times (68\times 4)\\&\color{answerred}=(25\times4)\times68\\&\color{answerred}=100\times68\\&\color{answerred}=6800\end{aligned}}
\MethodBlock{方法二：结合因数，先凑整}{用乘法交换律、结合律，先算容易凑整的两个因数。}{\begin{aligned}[t]&\phantom{={}}25\times (68\times 4)\\
&=(25\times 68)\times 4\\
&=1700\times 4\\
&=6800\end{aligned}}

\MethodBlock{方法三：先处理括号或乘除，再把特殊因数改成整百或整千}{先把25、125等改写成整百或整千除以一个数。}{\begin{aligned}[t]&\phantom{={}}25\times (68\times 4)\\
&=25\times 272\\
&=(100\div 4)\times 272\\
&=100\times (272\div 4)\\
&=100\times 68\\
&=6800\end{aligned}}

\MethodBlock{方法四：先处理括号或乘除，再凑整后补偿}{把接近整十、整百的因数凑整，再补回或减去多算的部分。}{\begin{aligned}[t]&\phantom{={}}25\times (68\times 4)\\
&=25\times 272\\
&=25\times (270+2)\\
&=25\times 270+25\times 2\\
&=6750+50\\
&=6800\end{aligned}}

\CalcBlock[方法一：简便计算]{9}{\begin{aligned}[t]&\phantom{={}}\color{black}65\times 37+65\times 62+65\\&\color{answerred}=65\times37+65\times62+65\times1\\&\color{answerred}=65\times(37+62+1)\\&\color{answerred}=65\times(99+1)\\&\color{answerred}=65\times100\\&\color{answerred}=6500\end{aligned}}
\MethodBlock{方法二：先处理括号或乘除，再交换、结合后先凑整}{保留每个数原来的加减号，先把能凑整的部分放在一起。}{\begin{aligned}[t]&\phantom{={}}65\times 37+65\times 62+65\\
&=2405+4030+65\\
&=(2405+65)+4030\\
&=2470+4030\\
&=6500\end{aligned}}

\MethodBlock{方法三：按运算顺序直接计算（对照）}{本法只作结果核对；本题要求简便计算，前面的简便解法仍保留。}{\begin{aligned}[t]&\phantom{={}}65\times 37+65\times 62+65\\
&=2405+4030+65\\
&=6435+65\\
&=6500\end{aligned}}

\par\noindent

\QHeading{第3题　解决问题。}

\Problem{水果店购进苹果、梨和香蕉各186箱，每箱苹果32元，每箱梨37元，每箱香蕉31元。水果店购进这些水果一共花了多少钱?}
\par


\MethodTitle{方法一}
\Answer{\Formula{\begin{aligned}[t]&\phantom{={}}(32+37+31)\times186\\&=(69+31)\times186\\&=100\times186\\&=18600\text{（元）}\end{aligned}}}
\Answer{答：一共花了18600元。}

\MethodTitle{方法二：分别计算三种水果的价钱}
\Answer{苹果、梨、香蕉的箱数相同，也可以先分别求总价再相加。}
\Answer{\Formula{\begin{aligned}[t]&\phantom{={}}32\times 186+37\times 186+31\times 186\\
&=5952+6882+5766\\
&=12834+5766\\
&=18600\text{（元）}\end{aligned}}}
\Answer{答：一共花了18600元。}
